\documentclass[10pt,a4paper]{amsart}
\usepackage[margin=19mm]{geometry}
\usepackage{amsmath,amssymb,mathtools}
\usepackage[scr=rsfs,cal=euler]{mathalfa}
\usepackage{xfrac}
\usepackage[unicode,hidelinks,bookmarks=false]{hyperref}
\newcommand{\ie}{i.e.\ }
\newcommand{\cf}{cf.\ }
\newcommand{\resp}{respectively\ }
\newcommand{\Subsection}{\subsection}
\providecommand{\nobreakdash}{\nobreak\hbox{-}}
\setlength{\emergencystretch}{2em}
\multlinegap0pt
\usepackage{siunitx}
\usepackage{siunitx}
\usepackage{siunitx}
\usepackage{siunitx}
\usepackage{mathtools}
\usepackage{amssymb}
\usepackage{tcolorbox}
\usepackage{enumerate}
\usepackage[shortlabels]{enumitem}
\usepackage{float}
\usepackage{algorithm}\usepackage{algpseudocode}
\usepackage{siunitx}
\newtheorem{lemma}{Lemma}
\newcommand{\Halmos}{\ensuremath{\quad\square}}
\title{Optimal Service Commitments in Traveling Salesman Problems with Stochastic Demand}
\author{Lina Schmidt, Julian Golak, Arne Schulz, Malte Fliedner}
\date{Source: arXiv:2609.11950v1 (CC BY 4.0)}
\begin{document}
\maketitle

\noindent\emph{Source and licence.} Optimal Service Commitments in Traveling Salesman Problems with Stochastic Demand. Version arXiv:2609.11950v1 (CC BY 4.0), distributed by arXiv under the Creative Commons Attribution 4.0 International licence, http://creativecommons.org/licenses/by/4.0/, as recorded in the arXiv licence metadata for this exact version and shown on the abstract page https://arxiv.org/abs/2609.11950v1. Full licence notice: \texttt{licenses/arxiv-2609.11950-CC-BY-4.0.txt}.\par\smallskip

\noindent\textit{Editorial scope.} Counted unit: the article's lemma lemma:discretisation\_length (source0.tex lines 291--296). The article's complete original proof of it is delivered with it (source0.tex lines 297--311, one \texttt{proof} environment of 15 lines). No mathematics is written, repaired or re-derived: the statement and the proof are the article's own lines, in the article's own notation, and no retained line is substituted or re-flowed.  Retained uncounted as the source-local context the proof needs: lines 234--236; lines 238--240.  Nothing was omitted from any retained span, so the package carries no omission record.  No retained line contains a citation, so the wrapper carries no bibliography and no bibliography entry.  No retained line contains a figure environment, an image-inclusion command or drawing code, so no image is delivered and the package declares no asset dependency.  Editorial formatting only, in addition: an AMS \texttt{amsart} preamble that restates the article's own declarations and macros used by the retained lines, namely the environments \texttt{lemma} on separate counters, together with the standard packages those lines need.  The retained environments print in this document as Lemma 1 in order of appearance, whereas the article prints the same items with the numbering of its own full document.  The complete native source of the article as distributed in the arXiv e-print for this exact version is bundled unchanged as \texttt{sources/210064/source0.tex}.
	\begin{equation}\label{eq:1stageOpt}
		\sup_{x \in \mathbb{R}_+} \mathbb{E}_{F\mid x}\left[ \text{rev}(x, F)\right].
	\end{equation}

	\begin{equation}\label{eq:2stageOpt}
		\text{rev}(x, F) = \max \left\{ \sum_{a \in S} r_a \colon S \subseteq F,\; \tau(S) \leq x \right\},
	\end{equation}

	\begin{lemma} \label{lemma:discretisation_length}
		The supremum in \eqref{eq:1stageOpt} is attained at a tour length, i.e.,
		\[
			\sup_{x \in \mathbb{R}_+} \mathbb{E}_{F\mid x}\left[ \text{rev}(x, F) \right] = \max_{i \in [K]} \, \mathbb{E}_{F\mid L_i}\left[ \text{rev}(L_i, F) \right].
		\]
	\end{lemma}

	\begin{proof}{Proof.}
		We first show that $\mathbb{E}_{F\mid x}\left[ \text{rev}(x, F) \right]$ is non-increasing on each interval $[L_i, L_{i+1})$, $i \in [K]$. Fix $i \in [K]$ and consider any $x_1, x_2 \in [L_i, L_{i+1})$ with $x_1 \leq x_2$. We observe that $\text{rev}(x_1, F) = \text{rev}(x_2, F)$ for all $F \subseteq A$: since $\mathcal{L}$ contains all tour lengths, no tour length lies in the open interval $(L_i, L_{i+1})$, so a subset $S \subseteq A$ satisfies $\tau(S) \leq x_1$ if and only if $\tau(S) \leq x_2$, both conditions being equivalent to $\tau(S) \leq L_i$. Hence the feasible sets in \eqref{eq:2stageOpt} coincide for $x_1$ and $x_2$, and we write $\text{rev}(F)$ for the common value. Moreover, $\text{rev}$ is monotone with respect to set inclusion: if $F \subseteq F'$, then every $S \subseteq F$ with $\tau(S) \leq L_i$ is also a feasible subset of $F'$, so $\text{rev}(F) \leq \text{rev}(F')$.

		Let $(U_a)_{a \in A}$ be independent random variables, each uniformly distributed on $[0,1]$, and define the random sets $F(x_j) \coloneqq \{a \in A \colon U_a \leq p_a(x_j)\}$ for $j \in \{1,2\}$. By independence,
		\begin{equation}\label{eq:help1}
			\mathbb{P}\left(F(x_j) = F\right) = \prod_{a \in F} p_a(x_j) \prod_{a \in A \setminus F} (1 - p_a(x_j)) = \mathbb{P}(F \mid x_j),
		\end{equation}
		for all \(F \subseteq A\),
		so $F(x_j)$ is distributed according to $\mathbb{P}(\,\cdot \mid x_j)$. Since each $p_a$ is non-increasing, $U_a \leq p_a(x_2)$ implies $U_a \leq p_a(x_1)$, and hence $F(x_2) \subseteq F(x_1)$ holds with probability one. Thus, we have $\text{rev}(F(x_2)) \leq \text{rev}(F(x_1))$ pointwise, and taking expectations gives
		\[
			\mathbb{E}_{F\mid x_2}\left[ \text{rev}(x_2, F) \right] = \mathbb{E}\left[ \text{rev}(F(x_2)) \right] \leq \mathbb{E}\left[ \text{rev}(F(x_1)) \right] = \mathbb{E}_{F\mid x_1}\left[ \text{rev}(x_1, F) \right].
		\]
		Hence, the expected revenue is non-increasing on $[L_i, L_{i+1})$ for every $i \in [K]$. Since these intervals partition $\mathbb{R}_+$, the supremum over each of them is attained at its left endpoint $L_i$, and the claim follows.
	\Halmos
\end{proof}

\end{document}
